\section{Gauge normal form and filtered residue adjoints}
\label{app:gauge-filtered-residues}

The normal-layer operator of \cref{sec:normal-complex} has a simpler
intrinsic form.  This section separates field-level solvability from the
finite Newton-window obstructions used in computation.

Write
\[
A=A_0,\qquad B=B_0,\qquad S=\alpha+\beta,
\]
and define the boundary one-form
\[
\Omega=\alpha A\,dB-\beta B\,dA.
\]
For a layer variation \((a,b)\), put
\[
u=\frac aA,\qquad v=\frac bB,\qquad
\xi=u+v,\qquad \eta=\alpha v-\beta u.
\]

\begin{proposition}[Gauge normal form]
\label{prop:gauge-normal-form}
The universal layer operator satisfies
\[
\mathscr D_r(a,b)
=AB\,\nabla_r\eta+\frac{S-r}{S}\xi\Omega,
\qquad
\nabla_r=d+\frac rS\,d\log(AB).
\]
If \(r\ne S\) and \(\Omega\ne0\), the operator is surjective over the
function field.  At the resonant layer \(r=S\),
\[
\mathscr D_S(a,b)
=d(\alpha Ab-\beta Ba).
\]
\end{proposition}

\begin{proof}
Substitute
\[
u=\frac{\alpha\xi-\eta}{S},\qquad
v=\frac{\beta\xi+\eta}{S}
\]
into the coordinate formula for \(\mathscr D_r\) and collect the
\(\xi\)- and \(\eta\)-terms.  Off resonance, for a rational differential
\(\Phi\), set \(\eta=0\) and
\[
\xi=\frac{S}{S-r}\frac{\Phi}{\Omega}.
\]
At \(r=S\), the multiplier term vanishes and
\(AB\nabla_S\eta=d(AB\eta)\).
\end{proof}

Thus, away from resonance, a cokernel can only be created by the finite pole
and Newton windows imposed on \(a\), \(b\), and the target differential.  It
is not a local obstruction to solving the rational differential equation.

\begin{proposition}[Extended deformation complex]
\label{prop:extended-complex}
Allow a variation \(\rho\Omega\) of the normalized right-hand side and put
\[
\widetilde{\mathscr D}_r(a,b,\rho)
=\mathscr D_r(a,b)-\rho\Omega.
\]
Then
\[
G_r(h)=\bigl(\alpha Ah,\beta Bh,(S-r)h\bigr)
\]
satisfies
\[
\widetilde{\mathscr D}_r\circ G_r=0.
\]
Consequently the normal layer carries a complex
\[
\mathcal G_r\xrightarrow{G_r}
\mathcal U_{A,r}\oplus\mathcal U_{B,r}\oplus\mathcal R_r
\xrightarrow{\widetilde{\mathscr D}_r}\mathcal W_r,
\]
whose degree \(-1,0,1\) terms represent infinitesimal gauge,
deformations modulo gauge, and obstructions.
\end{proposition}

\begin{proof}
The normal form gives
\[
\mathscr D_r(\alpha Ah,\beta Bh)=(S-r)h\Omega.
\]
\end{proof}

The same normal form clarifies the adjoint.  If \(M=AB\), integration by
parts gives
\[
\lambda\mathscr D_r(a,b)
=d(\lambda M\eta)
-M\eta\,\nabla_{S-r}\lambda
+\frac{S-r}{S}\lambda\xi\Omega.
\]
Thus
\[
(M\nabla_r)^\dagger=-M\nabla_{S-r}.
\]
The complementary-layer symmetry \(r\leftrightarrow S-r\) is invisible in
the original four-term coordinate formula.

\begin{theorem}[Filtered residue adjoint]
\label{thm:filtered-residue-adjoint}
Let \(U_{A,r},U_{B,r}\subset k(E)\) be the permitted correction windows and
let \(W_r\) be the permitted target window.  Represent \(W_r^\vee\) by
principal parts \(\lambda=(\lambda_p)\) at the allowed poles, with pairing
\[
\langle\lambda,\omega\rangle
=\sum_p\Res_p(\lambda_p\omega).
\]
Then
\[
(\operatorname{coker}\mathscr D_r)^\vee
\cong\ker\mathscr D_{r,\mathrm{win}}^\vee,
\]
where the filtered adjoint condition is that
\[
\beta B\,d\lambda+(S-r)\lambda\,dB
\]
annihilate \(U_{A,r}\), and
\[
-\alpha A\,d\lambda+(r-S)\lambda\,dA
\]
annihilate \(U_{B,r}\).  The equation
\[
\mathscr D_r(a_r,b_r)=-\Phi_r
\]
is solvable exactly when
\[
\sum_p\Res_p(\lambda_p\Phi_r)=0
\]
for every filtered adjoint class.
\end{theorem}

\begin{proof}
Multiply the coordinate operator by \(\lambda\), integrate each derivative
term by parts, and apply the residue theorem.  Boundary terms are precisely
the displayed principal-part pairings.  Finite-dimensional duality then
identifies the cokernel dual with the filtered adjoint kernel.
\end{proof}

There is an exact coefficient dictionary.  For target basis
\[
dz,z\,dz,\ldots,z^d\,dz
\]
and a left-null vector \(c=(c_0,\ldots,c_d)\), set
\[
\lambda_c(z)=\sum_{j=0}^dc_jz^{-j-1}.
\]
Then
\[
\Res_{z=0}\left(
\lambda_c\sum_{j=0}^df_jz^j\,dz
\right)=\sum_{j=0}^dc_jf_j.
\]
Thus a matrix null vector and its principal-part representative are exactly
the same functional in two coordinate systems.

\begin{corollary}[Resonant de Rham obstruction]
\label{cor:resonant-de-rham}
On \(E=\PP^1\), let \(\Sigma\) be the pole set and suppose the allowed
primitive window is complete modulo constants.  At \(r=S\),
\[
\operatorname{coker}\mathscr D_S
\cong
H^1_{\mathrm{dR}}(\PP^1\setminus\Sigma)
\cong
\left\{(\rho_p)_{p\in\Sigma}:\sum_p\rho_p=0\right\}.
\]
In particular, its dimension is \(|\Sigma|-1\).
\end{corollary}

\begin{proof}
At resonance the image consists of exact differentials.  A rational
differential on \(\PP^1\) has a rational primitive exactly when all residues
vanish, and the residue theorem supplies the single relation among them.
\end{proof}

\paragraph{Status of the six terminal components.}
The preceding results prove the operator normal form, gauge complex, and
filtered residue realization abstractly.
\Cref{prop:828-layer-injectivity} now proves that the ordinary full-support
layer maps have cokernel dimensions \(r-1\) for \(5\le r\le11\).
The exact artifact replay goes further in the original layer coordinates:
it reconstructs the stored matrices from
\(\mathscr D_r^{2,3}\), checks every filtered adjoint equation, and writes
every compatibility functional as a residue pairing with an explicit
principal part.  The nonzero pairing counts at layers \(6,7,8\) are
\(4,5,6\).  After the recorded normalization and duplicate removal, the
fifteen distinct equations have layer counts \(1,3,5,6\) at
\(r=5,6,7,8\), and they agree coefficientwise with the archived list.
The six equations used in the terminal contradiction are an explicit
coordinate projection of this fifteen-component residue section.

The remaining comparison is narrower: the archived normalized equations
have not been conjugated term by term into the new \((H,W,T)\)-coordinates.
Consequently we do not claim that the exact normal-coordinate change
preserves the required sparse Newton windows or identifies the admissible
approximate-root orbit.
